\documentclass{article}
\usepackage{eason}

\usepackage{pdfpages}

\usepackage[a4paper, margin=1in]{geometry}

\usepackage{lipsum}
\usepackage{multicol}

\newcommand{\DATE}{dummyDate}
\newcommand{\COMMITSHA}{dummySHA}

\title{TMUA Lecture 5 -- Graphs of Functions and Miscellaneous Questions}
\author{Eason Shao}
\date{v.\DATE.\COMMITSHA}

\begin{document}
\maketitle

\tableofcontents

\section{Functions and Graphs}
\subsection{Definitions}

For two sets \(A\) and \(B\), a \textbf{function}
\[
    \function{f}{A}{B}{x}{f\pqty{x}}
\]
is a relation where it maps \textbf{every} value \(a \in A\) to some \textbf{unique} value in \(b \in B\).

\begin{problem}
    How do we formalise this definition of a function?
\end{problem}

The set \(A\) is called its \textbf{domain} and the set \(B\) is called its \textbf{co-domain} (or range). In TMUA, we always have \(B = \RR\) and \(A\) being a nice subset of \(\RR\) (i.e., a union of intervals).

The set
\[
    \set{y \in B}{\exists x \in A: f\pqty{x} = y} = \set{f\pqty{x}}{x \in A}
\]
is the \textbf{image} or \textbf{range} of \(f\) (note the dual-use of the word range with different meanings). It is the set of values in the co-domain that the function can take.

The \textbf{vertical line test} tests whether some graph (of a relation) is a function or not.

\begin{example}
    Determine if \(x = y^3\) and \(x = y^2\) are functions.
\end{example}

\begin{problem}
    What is the point of having the co-domain in the definition?
\end{problem}

A function is \textbf{one-to-one}, or \textbf{injective}, if for all \(a \neq b\) in its domain, we have
\[
    f\pqty{a} \neq f\pqty{b}.
\]

Otherwise, we say a function is \textbf{many-to-one}.

The \textbf{horizontal line test} tests whether some graph (of a function) is one-to-one or many-to-one.

\begin{example}
    Determine if \(y = x^2\) and \(y = x\) are one-to-one or many-to-one.
\end{example}

\subsection{Properties of Functions}
\subsubsection{Monotonicity}

A function \(f\) is \textbf{increasing} on a interval \(\ccintv{a}{b}\) (which must be entirely inside its domain), if for all \(a \leq x_1 \leq x_2 \leq b\), we have
\[
    f\pqty{x_1} \leq f\pqty{x_2}.
\]

Changing \(f\pqty{x_1} \geq f\pqty{x_2}\) gives the definition of a \textbf{decreasing} function.

Recall from definition that we may use \(f'\) to see whether a function is increasing or decreasing.

\subsubsection{Parity}

If a function's domain is symmetric about the origin, we may discuss the \textbf{parity} of a function.

A function \(f\) is \textbf{odd} if
\[
    f\pqty{x} = - f\pqty{-x},
\]
and is \textbf{even} if
\[
    f\pqty{x} = f\pqty{-x}.
\]

The words odd and even come from the powers of \(x\) -- the function \(x\) raised to an even/odd power is even/odd.

\begin{remark}
    You may sometimes hear people refer to this as \textit{``odevity''} (in China) -- this is \textbf{not a proper terminology}.
\end{remark}

The most important application of odd/even functions is in integration. If \(f\) is odd, then
\[
    \int_{-a}^{a} f\pqty{x} \dd{x} = 0
\]
and if \(f\) is even, then
\[
    \int_{-a}^{a} f\pqty{x} \dd{x} = 2 \int_{0}^{a} f\pqty{x} \dd{x}.
\]

Note a function can be neither even nor odd, or both even and odd.

\begin{example}
    If \(f\) is a function defined on \(\RR\) and is both even and odd, what is \(f\)?
\end{example}

\subsubsection{Periodicity}

For a function \(f\) defined on \(\RR\), we say some real number \(T \neq 0\) is a \textbf{period} of \(f\), if
\[
    f \pqty{x + T} = f \pqty{x}
\]
for all \(x\). If such \(T\) exists, we say \(f\) is \textbf{periodic}.

If we refer to \textbf{the period} of a function, we mean the \textbf{least positive period} (which does not necessarily exist).

For example, the \(\sin\) function has all \(2 k \pi\) for \(k \neq 0\) and \(k \in \ZZ\) as its period, but \textbf{the period} of \(\sin\) is \(2 \pi\).

\begin{problem}
    A constant function is periodic but it does not have the least positive period. Can you find a non-constant function with this property?
\end{problem}

\subsubsection{Symmetry}

We say a function \(f\) is \textbf{symmetric} about the line \(x = a\), if
\[
    f\pqty{a + x} = f\pqty{a - x}
\]  
for all \(x\). Equivalently, this could be written as
\[
    f\pqty{x} = f\pqty{2a - x}
\]
for all \(x\).

This gives the integration property
\[
    \int_{a - x}^{a + x} f\pqty{x} \dd{x} = 2 \int_{a}^{a + x} f\pqty{x} \dd{x}.
\]

A function \(f\) is \textbf{rotationally symmetric by \(180 \degree\)} about a point \(\pqty{a, f\pqty{a}}\), if
\[
    f\pqty{a + x} - f\pqty{a} = f\pqty{a} - f\pqty{a - x}
\]
for all \(x\). Equivalently, this could be written as
\[
    f\pqty{a + x} + f\pqty{a - x} = 2 f\pqty{a},
\]
or
\[
    f\pqty{x} + f\pqty{2a - x} = 2 f\pqty{a}
\]
for all \(x\).

This gives the integration property
\[
    \int_{a - x}^{a + x} f\pqty{x} \dd{x} = 2 f\pqty{a} x.
\]

\textbf{Note parity is a special case of symmetry where \(a = 0\).}

\begin{example}
    Integrate
    \[
        \int_{0}^{4} x \pqty{x - 1} \pqty{x - 2} \pqty{x - 3} \pqty{x - 4} \dd{x}.
    \]
\end{example}

\subsection{Graphs of Common Functions}

\textbf{Linear, Quadratic, Polynomials, Trigonometric} -- see previous lectures.

The \textbf{absolute value} or \textbf{modulus function} \(y = \abs{x}\) is defined as
\[
    y = \abs{x} = \begin{cases}
        x & \qif x \geq 0, \\
        -x & \qif x < 0.
    \end{cases}
\]

The \textbf{square root function} \(y = \sqrt{x}\) is \textbf{only defined on} \(x \geq 0\) and has range \(y \geq 0\).

Let \(a > 0\). \textbf{Exponential functions} are functions of the form
\[
    f\pqty{x} = a^x \qqtext{where} a > 0,
\]
and \textbf{logarithmic functions} are their inverses, \(f\inverse\pqty{x} = \log_a x\) for \(a \neq 1\).

The graphs differ for the cases \(0 < a < 1\) (in which case both are decreasing), \(a = 1\) (in which the logarithm does not exist) and \(a > 1\) (in which case both are increasing) -- \textbf{be very careful of the differing monotonic properties}.

\subsection{Graph Transformations}

\begin{itemize}
    \item \textbf{The graph of \(\abs{f\pqty{x}}\)} is formed by:
        \begin{itemize}
            \item sketching \(y = f\pqty{x}\); then
            \item reflecting the graph below the \(x\)-axis (\(y < 0\)) in the \(x\)-axis, removing the original.
        \end{itemize}

    \item \textbf{The graph of \(f\pqty{\abs{x}}\)} is formed by:
        \begin{itemize}
            \item sketching \(y = f\pqty{x}\) for \(x \geq 0\); then
            \item reflecting the graph to the right of \(y\)-axis (\(x > 0\)) in the \(y\)-axis, keeping the original.
        \end{itemize}

    \item \textbf{The graph of \(f\pqty{x + a}\)} is formed by translating the graph of \(y = f\pqty{x}\) by \(a\) units in the \textbf{negative \(x\)-direction}.
    \item \textbf{The graph of \(f\pqty{x} + a\)} is formed by translating the graph of \(y = f\pqty{x}\) by \(a\) units in the \textbf{positive \(y\)-direction}.
    \item \textbf{The graph of \(f\pqty{ax}\)} is formed by stretching the graph of \(y = f\pqty{x}\) \textbf{in the \(x\)-axis by a scale factor of \(\flatfrac{1}{a}\)}.
    \item \textbf{The graph of \(a f\pqty{x}\)} is formed by stretching the graph of \(y = f\pqty{x}\) \textbf{in the \(y\)-axis by a scale factor of \(a\)}.
        
\end{itemize}

\begin{example}
    Sketch the following functions:
    \begin{multicols}{3}
        \begin{itemize}
            \item \(\abs{x^2 - 4}\)
            \item \(\abs{\sin x}\)
            \item \(3 - \abs{x}\)
            \item \(e^{\abs{x}}\)
            \item \(2^{x + 1}\)
            \item \(\sin\pqty{x + \frac{\pi}{6}}\)
            \item \(x^3 + 3\)
            \item \(\cos x - 2\)
            \item \(\sqrt{2x}\)
            \item \(\sin \pqty{3x}\)
            \item \(5 \cos x\)
            \item \(- \frac{1}{2} \sqrt{x}\)
        \end{itemize}
    \end{multicols}
\end{example}

\begin{example}[TMUA 2017 P1 Q18]
    The graph of \(y = \log_{10} x\) is translated in the positive \(y\)-direction by \(2\) units. This translation is equivalent to a stretch of factor \(k\) parallel to the \(x\)-axis. Find \(k\).
\end{example}

\begin{example}[MAT 2014 Question I]
    The graph of the function
    \[
        y = 2^{x^2 - 4x + 3}
    \]
    can be obtained from the graph of \(y = 2^{x^2}\) by
    \begin{enumerate}[label=(\alph*)]
        \item a stretch parallel to the \(y\)-axis followed by a translation parallel to the \(y\)-axis.
        \item a translation parallel to the \(x\)-axis followed by a stretch parallel to the \(y\)-axis.
        \item a translation parallel to the \(x\)-axis followed by a stretch parallel to the \(x\)-axis.
        \item a translation parallel to the \(x\)-axis followed by reflection in the \(y\)-axis.
        \item reflection in the \(y\)-axis followed by translation parallel to the \(y\)-axis.
    \end{enumerate}
\end{example}

\subsection{Application of Graphs of Functions}

Previously, we have seen the application of using graphs of functions to solve (qualitatively) a system of inequalities.

We may also use graphs to solve a system of equations: the intersections of graphs correspond to solutions to the original system of equations.

Similarly, the intersections of the graphs of \(f\) and \(g\) solve the equation \(f\pqty{x} = g\pqty{x}\). (Sometimes, we re-arrange the functions \(f\) and \(g\) to make the graph-sketching easier.)

\begin{example}[TMUA 2016 P1 Q10]
    How many solutions does the equation \(x \tan x = 1\) have in the interval \(-2 \pi \leq x \leq 2\pi\)?
\end{example}

\section{Miscellaneous Questions}

\subsection{Ratios}

Key idea: set the total to be \(1\) (or \(10\), or \(100\)).

\begin{example}[TMUA 2016 P1 Q6]
    \(P\), \(Q\) and \(R\) are each mixtures of red and white paint. The percentage by volume of red paint in \(P\) is \(\SI{30}{\percent}\). The percentage by volume of red paint in \(Q\) is \(\SI{20}{\percent}\).

    The mixtures \(P\), \(Q\) and \(R\) are combined in the proportion \(12 : 5 : 3\) respectively.

    If the resulting mixture contains \(\SI{25}{\percent}\) by volume of red paint, what percentage by volume of mixture \(R\) is red paint?
\end{example}

\begin{example}[TMUA 2016 P1 Q7]
    \(\SI{60}{\percent}\) of a sports club's members are women and the remainder are men. This sports club offers the opportunity to play tennis or cricket. Every member plays exactly one of the two sports.

    \(\flatfrac{2}{5}\) of the male members of the club play cricket; \(\flatfrac{2}{3}\) of the cricketing members of the club are women.

    What is the probability that a member of the club, chosen at random, is a woman who plays tennis?
\end{example}

\subsection{Logical Deduction}

\begin{example}[TMUA 2017 P2 Q20]
    I have forgotten my \(5\)-character computer password, but I know that it consists of the letters \texttt{a}, \texttt{b}, \texttt{c}, \texttt{d}, \texttt{e} in some order. When I enter a potential password into the computer, it tells me exactly how many of the letters are in the correct position.

    When I enter \texttt{abcde}, it tells me that none of the letters are in the correct position. The same happens when I enter \texttt{cdbea} and \texttt{eadbc}.

    Using the best strategy, how many \textbf{further} attempts must I make in order to \textbf{guarantee} that I can \textbf{deduce} the correct password?
\end{example}

\end{document}